% !TEX encoding = UTF-8 Unicode
% !TEX root = thesis.tex

\chapter*{Summary and Conclusions}

This thesis investigated machine-learning-based, drift-aware monitoring for
handover (HO) optimization in 5G/6G radio access networks. Because no public
5G-SA dataset exposes both the HO control parameters and the resulting mobility
KPIs, the main data was produced with a custom, 3GPP-compliant 5G NR SA
handover simulator; its radio layer was calibrated against the real NordicDat
drive-test dataset and additionally checked through an external face-validity
test, with the real data never merged into the simulated benchmark. Returning to
the aim and the four research questions:

\begin{itemize}
  \item \textbf{RQ1 (anomaly detection):} unsupervised detectors catch
  structured faults well (measurement-glitch PR-AUC up to $0.858$, RLF-burst up
  to $0.751$) but subtle faults (interference spikes and slow single-UE
  degradation) stay near the noise floor; pooled over the timeline Isolation
  Forest is the strongest single detector (PR-AUC $0.425$), and an ablation shows
  the RSRQ feature family is the most informative by a wide margin.
  \item \textbf{RQ2 (drift):} every injected drift is detected by every detector
  (no misses), so the practical question is a latency-versus-false-alarm
  trade-off---batch two-sample tests (Energy/MMD) fire in $\approx 3$\,s at a
  high false-alarm rate, whereas sequential detectors (ADWIN/DDM) keep false
  alarms near $0.005$ per second but react more slowly.
  \item \textbf{RQ3 (adaptation):} the central result, and a somewhat surprising one,
  is that \emph{how} a monitor is retrained matters more than \emph{whether} it
  is retrained. Naive periodic or drift-triggered refits roughly halve PR-AUC
  ($0.264 \rightarrow \approx 0.12$) through self-poisoning, whereas a
  self-supervised filtered refit preserves the static baseline overall
  ($0.248$ vs $0.264$) and more than triples drift-phase PR-AUC
  ($0.057 \rightarrow 0.187$), at lower cost.
  \item \textbf{RQ4 (performance impact and reconfiguration):} a
  gradient-boosting configuration--performance surrogate predicts mobility KPIs
  accurately (HOSR $R^2 = 0.96$, RLF $R^2 = 0.95$, ping-pong $R^2 = 0.72$) with
  split-conformal intervals that restore calibrated coverage; an inverse query
  returns $12$ of $36$ feasible configurations (favouring the shortest
  time-to-trigger), and an end-to-end detect--recommend--recover loop fires seven
  interventions for a cumulative, model-predicted HOSR uplift of $+0.62$.
\end{itemize}

\section*{Limitations}

The scope is 5G NR SA intra-RAT inter-gNB Xn handover; NSA, intra-gNB and
inter-RAT mobility are out of scope. The framework is deliberately positioned as
monitoring and decision support rather than as a parameter-tuning control loop:
automatic tuning of time-to-trigger, hysteresis and cell individual offset is
already standardized as mobility robustness optimization (MRO)~\cite{tr36839,ts28313},
and the reconfiguration step here is intended to feed such a function rather than
to replace it. In contrast to MRO, which acts on lagging mobility-failure
counters, the system intervenes on leading drift indicators; it does not,
however, forecast the failure of an individual handover before it occurs. The
anomalies are \emph{injected} in
simulation rather than observed in the field; the reported HOSR uplift is a
surrogate \emph{model prediction}, not a measured field gain; the RSRQ radio
layer retains a residual calibration mismatch (KS $0.467$); and the
per-anomaly-type best detector is seed-sensitive, although the pooled detector
ranking and the filtered-over-naive result are stable across seeds and scenarios.

\section*{Future work}

Several concrete directions follow from the limitations above. The most immediate
is to bring the experimental environment closer to reality without relying on
operator event logs, which are rarely accessible in practice. This means richer
system-level simulators (for example ns-3 5G-LENA or Simu5G) and, longer term,
digital-twin style environments, while keeping the controlled injection of drift
and anomalies that makes labelled benchmarking possible; public drive-test
datasets such as NordicDat and the urban Bangladesh LTE handover set can continue
to serve as external face-validity references.

A second direction is continual (online, incremental) learning, so that the
anomaly monitor and the configuration-performance surrogate are updated
continuously under persistent change rather than refitted in batches. The
self-supervised filtering rule introduced here is a natural starting point and
should be made adaptive and given stability guarantees. A third direction is
explainability: identifying which KPIs, radio conditions, or parameter changes
are most strongly associated with handover degradation, so that a recommended
reconfiguration can be justified rather than only predicted. A related direction
is to move from detection toward genuine \emph{proactive} prediction: the present
framework reacts to the onset of drift on leading KPI streams, but it does not
forecast the failure of an individual handover before it happens. Adding an
event-level predictor (for example a per-UE time-to-failure or
handover-failure model that estimates the probability of a too-late or too-early
handover ahead of time) was kept out of the present scope because it requires
per-event prediction targets and a forecasting setup distinct from the
unsupervised monitoring studied here; combined with the configuration--performance
surrogate, such a predictor would let the system act before the precursor drift
fully materialises and connects naturally to the proactive handover-prediction
line of work reviewed in Section~\ref{subsec:rw-prediction}. Finally, the
configuration scope can be broadened beyond time-to-trigger, hysteresis, and A3
offset to include cell individual offset (CIO), and beyond intra-RAT inter-gNB to
the cases excluded here, moving toward a safe and interpretable decision-support
mechanism that limits unnecessary optimisation actions. Together these steps
extend handover monitoring toward the broader goal of drift-aware, explainable,
and adaptive mobility management.
